\section{Исследование и построение решения задачи}
\label{sec:Chapter3} \index{Chapter3}

\subsection{Декомпозиция задачи}

Исходная задача восстановления речи разбивается на следующие подзадачи:
\begin{enumerate}
    \item выбор открытых SSL-энкодеров, способных заменить закрытый USM
        (подраздел~\ref{ssec:enc-choice});
    \item проектирование обучаемого очистителя признаков, встраиваемого в
        замороженный энкодер (подраздел~\ref{ssec:cleaner});
    \item выбор и адаптация нейросетевых вокодеров под выход очистителя
        (подраздел~\ref{ssec:voc-choice});
    \item определение слоя извлечения признаков, оптимального для восстановления
        (подраздел~\ref{ssec:layers});
    \item формирование данных и метрик для экспериментального сравнения
        (подразделы~\ref{ssec:data}--\ref{ssec:metrics}).
\end{enumerate}

\subsection{Архитектура предлагаемой системы}
\label{ssec:arch}

Предлагаемая система воспроизводит трёхкомпонентную схему Miipher-2 с открытыми
блоками (рис.~\ref{fig:arch}). Замороженный SSL-энкодер $E$ преобразует
искажённый сигнал $y$ в последовательность признаков слоя $\ell$:
$z^{(\ell)} = E_\ell(y)$. Очиститель $G_\phi$ предсказывает чистые признаки
$\hat{z}^{(\ell)}$, которые подаются на вокодер $V_\psi$, синтезирующий
$\hat{x}$. Энкодер заморожен, что резко сокращает число обучаемых параметров и
позволяет переиспользовать богатые предобученные представления.

\begin{figure}[H]
    \centering
    \begin{tikzpicture}[
        >={Stealth[length=2.4mm]},
        lyr/.style={draw, minimum width=3.4cm, minimum height=0.72cm,
            align=center, font=\scriptsize, fill=black!6},
        sel/.style={lyr, fill=black!16, very thick},
        adp/.style={draw, rounded corners=2pt, minimum width=1.3cm,
            minimum height=0.6cm, align=center, font=\scriptsize, fill=black!22},
        sum/.style={draw, circle, minimum size=0.42cm, inner sep=0pt,
            font=\scriptsize},
        io/.style={align=center, font=\scriptsize}]

        % --- вертикальный стек энкодера ---
        \node (cnn) [lyr] {Свёрточный фронтенд (CNN)};
        \node (l1)  [lyr, above=1.0cm of cnn] {Трансформерный слой $1$};
        \node (ld)  [above=0.55cm of l1, font=\large] {$\vdots$};
        \node (le)  [sel, above=0.55cm of ld] {Трансформерный слой $\ell$};

        % вход снизу
        \node (in) [io, below=0.75cm of cnn] {Искажённый\\сигнал $y$};
        \draw[->] (in) -- (cnn);
        \draw[->] (cnn) -- (l1);
        \draw[->] (l1) -- (ld);
        \draw[->] (ld) -- (le);

        % --- параллельные адаптеры справа от слоёв ---
        \node (a1) [adp, right=1.7cm of l1] {Адаптер};
        \node (ae) [adp, right=1.7cm of le] {Адаптер};
        \draw[->] (l1.east) -- (a1.west);
        \draw[->] (le.east) -- (ae.west);

        % --- пунктирная рамка энкодера+очистителя ---
        \begin{scope}[on background layer]
            \node (encbox) [draw, dashed, rounded corners,
                fit=(cnn)(le)(a1)(ae), inner sep=8pt] {};
        \end{scope}
        \node [font=\scriptsize, anchor=south] at (encbox.north)
            {Замороженный энкодер $E$ + обучаемые адаптеры $G_\phi$};

        % --- выход выбранного слоя в вокодер ---
        \node (voc) [draw, rounded corners, fill=black!16, minimum width=2.6cm,
            minimum height=1.0cm, align=center, font=\scriptsize,
            right=2.3cm of ae] {Вокодер $V_\psi$\\(HiFi-GAN / WaveFit)};
        \node (res) [io, below=1.0cm of voc] {Чистый\\сигнал $\hat{x}$};
        \draw[->] (ae.east) -- node[above, font=\tiny] {$\hat{z}^{(\ell)}$} (voc.west);
        \draw[->] (voc) -- (res);
    \end{tikzpicture}
    \caption{Архитектура системы: параллельные адаптеры встраиваются в
        замороженный энкодер, признаки слоя $\ell$ подаются на вокодер}
    \label{fig:arch}
\end{figure}

\subsection{Выбор открытых энкодеров}
\label{ssec:enc-choice}

В качестве замены USM выбраны две открытые модели с архитектурой,
сопоставимой по принципу работы:
\begin{itemize}
    \item \textbf{HuBERT Base} --- 12 трансформерных слоёв, обучение
        предсказанием кластеров; представляет контрастно-кластерное
        направление SSL;
    \item \textbf{WavLM Base+} --- 12 слоёв, обучение с моделированием шума и
        наложенной речи; благодаря денойзингу при предобучении ожидается большая
        устойчивость к деградациям.
\end{itemize}
Обе модели имеют скрытую размерность 768 и частоту кадров 50~Гц (шаг 20~мс), что
удобно для согласования с вокодером. Энкодеры используются в замороженном виде:
обучаются только адаптеры и вокодер.

\subsection{Очиститель признаков}
\label{ssec:cleaner}

Очиститель реализован в виде \emph{параллельных адаптеров} --- компактных
блоков прямого распространения, подключаемых к выходам трансформерных слоёв
энкодера. Каждый адаптер представляет собой сужающе-расширяющую пару линейных
слоёв с нелинейностью GELU и остаточной связью:
\begin{equation}
    \tilde{h}^{(i)} = h^{(i)} +
    W_{\text{up}}\,\sigma\!\big(W_{\text{down}} h^{(i)}\big),
    \label{eq:adapter}
\end{equation}
где $h^{(i)}$ --- выход $i$-го слоя энкодера, $W_{\text{down}}\in
\mathbb{R}^{256\times 768}$, $W_{\text{up}}\in\mathbb{R}^{768\times256}$,
$\sigma$ --- функция активации. Скрытая размерность адаптера 256 выбрана как
компромисс между ёмкостью и числом параметров. Суммарно адаптеры добавляют около
$4{,}9$~млн обучаемых параметров. Обучение ведётся по функции
потерь~\eqref{eq:fcloss} с $\lambda = 1$; целевые признаки $z^{\ast}$ получаются
прогоном чистого сигнала через тот же замороженный энкодер.

\subsection{Выбор вокодеров}
\label{ssec:voc-choice}

Рассмотрены два открытых вокодера, принимающих на вход последовательность
SSL-признаков:
\begin{itemize}
    \item \textbf{HiFi-GAN}~\cite{hifigan} (конфигурация V1) --- однопроходный
        GAN-вокодер; быстрый и компактный, входной проекционный слой адаптирован
        под размерность 768;
    \item \textbf{WaveFit}~\cite{wavefit} (5 итераций) --- итеративный вокодер,
        используемый в оригинальном Miipher; обеспечивает более естественное
        звучание ценой пятикратного увеличения вычислений.
\end{itemize}
Сочетание двух энкодеров и двух вокодеров даёт \textbf{четыре комбинации}
<<энкодер — вокодер>>, исследуемые далее (табл.~\ref{tab:configs}).

\begin{table}[H]
    \caption{Исследуемые конфигурации системы}
    \label{tab:configs}
    \centering
    \footnotesize
    \begin{tabularx}{\textwidth}{l l l X}
        \toprule
        Обозначение & Энкодер & Вокодер & Комментарий \\
        \midrule
        K1 & HuBERT Base & HiFi-GAN & базовая быстрая конфигурация \\
        K2 & HuBERT Base & WaveFit  & итеративный вокодер \\
        K3 & WavLM Base+ & HiFi-GAN & устойчивый энкодер + быстрый вокодер \\
        K4 & WavLM Base+ & WaveFit  & наиболее близка к Miipher-2 \\
        \bottomrule
    \end{tabularx}
\end{table}

\subsection{Функция потерь вокодера}
\label{ssec:vocloss}

Вокодер обучается в генеративно-состязательном режиме. Пусть $V_\psi$ ---
генератор, $D$ --- набор дискриминаторов (многомасштабный и многопериодный, как
в~\cite{hifigan}). Состязательная составляющая для генератора и дискриминатора
имеет вид
\begin{equation}
    \mathcal{L}_{\text{adv}}(V_\psi) = \mathbb{E}\big[(D(\hat{x})-1)^2\big],
    \qquad
    \mathcal{L}_{\text{adv}}(D) = \mathbb{E}\big[(D(x)-1)^2 + D(\hat{x})^2\big].
    \label{eq:advloss}
\end{equation}
Выражение~\eqref{eq:advloss} соответствует функционалу метода наименьших
квадратов (LS-GAN). К нему добавляются мел-спектральная потеря
$\mathcal{L}_{\text{mel}} = \|S(x)-S(\hat{x})\|_1$, вычисляемая по~\eqref{eq:mel},
и потеря согласования признаков дискриминатора $\mathcal{L}_{\text{fm}}$. Итоговая
функция потерь генератора --- взвешенная сумма
\begin{equation}
    \mathcal{L}_{V} = \mathcal{L}_{\text{adv}}(V_\psi)
    + \lambda_{\text{fm}}\mathcal{L}_{\text{fm}}
    + \lambda_{\text{mel}}\mathcal{L}_{\text{mel}},
    \label{eq:vtotal}
\end{equation}
где $\lambda_{\text{fm}}=2$, $\lambda_{\text{mel}}=45$ --- весовые коэффициенты
в~\eqref{eq:vtotal}, типичные для HiFi-GAN. Мел-спектральная потеря отвечает за точность
воспроизведения, состязательная --- за естественность и детализацию сигнала.

\subsection{Исследование слоя извлечения признаков}
\label{ssec:layers}

Поскольку слои SSL-энкодера кодируют разную информацию, выбор слоя $\ell$
является ключевым гиперпараметром. Для его подбора проведён предварительный
эксперимент: при фиксированном вокодере HiFi-GAN обучались адаптеры на признаках
каждого слоя $\ell\in\{1,\dots,12\}$, после чего на валидационном наборе
измерялся PESQ. Результаты для обоих энкодеров приведены на
рис.~\ref{fig:layers}.

\begin{figure}[H]
    \centering
    \begin{tikzpicture}
    \begin{axis}[
        width=0.82\textwidth, height=6.2cm,
        xlabel={Номер слоя $\ell$}, ylabel={PESQ (валидация)},
        xmin=0.5, xmax=12.5, ymin=2.2, ymax=3.1,
        xtick={1,2,3,4,5,6,7,8,9,10,11,12},
        legend pos=south west, legend cell align=left,
        grid=both, grid style={black!12}, tick label style={font=\scriptsize},
        label style={font=\small}]
    \addplot[mark=*, thick] coordinates {
        (1,2.39)(2,2.57)(3,2.65)(4,2.82)(5,2.84)(6,2.93)(7,2.95)(8,2.87)
        (9,2.85)(10,2.71)(11,2.69)(12,2.55)};
    \addlegendentry{WavLM Base+}
    \addplot[mark=square*, thick, dashed] coordinates {
        (1,2.36)(2,2.44)(3,2.61)(4,2.66)(5,2.79)(6,2.82)(7,2.78)(8,2.76)
        (9,2.64)(10,2.62)(11,2.50)(12,2.47)};
    \addlegendentry{HuBERT Base}
    \end{axis}
    \end{tikzpicture}
    \caption{Зависимость качества восстановления (PESQ) от номера слоя
        извлечения признаков}
    \label{fig:layers}
\end{figure}

Для обоих энкодеров наблюдается выраженный максимум в области \emph{средних}
слоёв: $\ell = 7$ для WavLM и $\ell = 6$ для HuBERT. Это согласуется с данными
SUPERB~\cite{superb}: средние слои содержат наиболее устойчивое фонетическое
представление, тогда как нижние слои сохраняют слишком много акустического шума,
а верхние --- слишком абстрактны и теряют детали, необходимые для качественного
ресинтеза. Примечательно, что выбранный для USM 13-й слой из 32~\cite{miipher2}
также относится к средней трети сети, что подтверждает переносимость наблюдения.
В дальнейших экспериментах используются слои $\ell=7$ (WavLM) и $\ell=6$
(HuBERT).

\subsection{Данные}
\label{ssec:data}

Для обучения и оценки сформирован набор пар <<чистый — искажённый>>. Чистая речь
взята из корпусов LibriTTS~\cite{libritts} (англоязычный, многоговорящий) и
VCTK~\cite{vctk}. Искажения синтезировались согласно
модели~\eqref{eq:degradation}: аддитивный шум из баз DEMAND~\cite{demand} и DNS
Challenge~\cite{dnschallenge} подмешивался с равномерно распределённым
$\mathrm{SNR}\in[0;20]$~дБ; реверберация моделировалась сгенерированными RIR;
часть примеров дополнительно сжималась кодеками. Применялись четыре схемы
деградации (шум; шум+реверберация; шум+кодек; все три), как и в~\cite{miipher2}.
Распределение данных приведено в табл.~\ref{tab:data}.

\begin{table}[H]
    \caption{Состав используемых данных}
    \label{tab:data}
    \centering
    \footnotesize
    \begin{tabularx}{\textwidth}{X c c c}
        \toprule
        Выборка & Источник чистой речи & Объём, ч & Назначение \\
        \midrule
        Обучающая    & LibriTTS train     & 240 & обучение адаптеров и вокодера \\
        Валидационная & LibriTTS dev       & 12  & подбор слоя и гиперпараметров \\
        Тестовая A   & VCTK (тест)         & 8   & оценка на невиданных дикторах \\
        Тестовая B   & LibriTTS test-clean & 10  & сравнение с известными корпусами \\
        \bottomrule
    \end{tabularx}
\end{table}

\subsection{Процедура аугментации}
\label{ssec:augment}

Зашумлённые примеры синтезировались динамически в процессе обучения, что
обеспечивает практически неограниченное разнообразие пар. Для чистого сигнала
$x$ и случайно выбранного шумового фрагмента $n$ смесь формировалась по
правилу~\eqref{eq:mix} с заданным отношением сигнал/шум:
\begin{equation}
    y_0 = (x * h) + g\cdot n,
    \qquad
    g = \sqrt{\frac{\|x*h\|^2}{\|n\|^2}\cdot 10^{-\mathrm{SNR}/10}},
    \label{eq:mix}
\end{equation}
где масштабирующий коэффициент $g$ подбирается так, чтобы достичь требуемого
значения $\mathrm{SNR}$ согласно~\eqref{eq:snr}. Импульсные характеристики
помещения $h$ генерировались методом источников-изображений для случайных
размеров комнаты и времён реверберации $T_{60}\in[0{,}2; 0{,}8]$~с. Для части
примеров к $y_0$ дополнительно применялось сжатие кодеками (MP3, OPUS, $\mu$-law)
с последующим декодированием, моделирующее артефакты передачи. Итоговое
распределение четырёх схем деградаций задавалось равновероятным.

\subsection{Представление сигнала и признаки}
\label{ssec:features}

Все сигналы приводились к моночастоте дискретизации 16~кГц. SSL-энкодер
принимает сырую волновую форму и через свёрточный фронтенд формирует
последовательность кадров с шагом 20~мс. Вокодер HiFi-GAN при обучении
дополнительно использует мел-спектрограмму $S\in\mathbb{R}^{F\times M}$,
вычисляемую как
\begin{equation}
    S(f,m) = \log\!\Big(\textstyle\sum_{k} \mathrm{Mel}_f(k)\,|Y(k,m)|^2 +
    \varepsilon\Big),
    \label{eq:mel}
\end{equation}
где $\mathrm{Mel}_f(k)$ --- $f$-й мел-фильтр, $\varepsilon$ --- константа
численной устойчивости. Параметры STFT: длина окна 1024 отсчёта, шаг 256
отсчётов, число мел-каналов $F=80$. Мел-спектрограмма используется в составе
функции потерь вокодера для согласования с целевым сигналом.

\subsection{Метрики оценки}
\label{ssec:metrics}

Качество восстановления оценивается набором объективных метрик.
\begin{itemize}
    \item \textbf{PESQ}~\cite{pesq} --- перцептивная оценка качества в диапазоне
        от $-0{,}5$ до $4{,}5$; интрузивная (требует эталона).
    \item \textbf{STOI}~\cite{stoi} --- оценка разборчивости в диапазоне
        $[0;1]$.
    \item \textbf{SI-SDR}~\cite{sisdr} --- масштабно-инвариантное отношение
        сигнал/искажение в дБ:
        \begin{equation}
            \mathrm{SI\text{-}SDR} = 10\lg
            \frac{\|\alpha x\|^2}{\|\alpha x - \hat{x}\|^2},
            \qquad
            \alpha = \frac{\langle \hat{x}, x\rangle}{\|x\|^2}.
            \label{eq:sisdr}
        \end{equation}
        Масштабный множитель $\alpha$ в~\eqref{eq:sisdr} обеспечивает
        инвариантность метрики к усилению сигнала.
    \item \textbf{DNSMOS}~\cite{dnsmos} --- неинтрузивная нейросетевая оценка
        (не требует эталона), коррелирующая с субъективным MOS.
\end{itemize}
Вычислительная эффективность характеризуется коэффициентом реального времени
(RTF) и числом обучаемых параметров. Совокупность метрик позволяет оценить как
качество, так и практическую применимость каждой конфигурации.
